% =====================================================================
%  preamble.tex  --  packages, page geometry, styles and book macros
%  Compile with XeLaTeX (see latexmkrc / Makefile).
% =====================================================================

% ---------- fonts -----------------------------------------------------
\usepackage{tocloft}
\setcounter{tocdepth}{2}
\renewcommand{\cftchapfont}{\sffamily\bfseries}
\renewcommand{\cftchappagefont}{\sffamily\bfseries}
\renewcommand{\cftsecfont}{\sffamily\small}
\renewcommand{\cftsecpagefont}{\sffamily\small}
\setlength{\cftbeforechapskip}{0.9em}

\usepackage{fontspec}
\usepackage{unicode-math}
\setmainfont{texgyrepagella}[
  Extension      = .otf,
  UprightFont    = *-regular,
  BoldFont       = *-bold,
  ItalicFont     = *-italic,
  BoldItalicFont = *-bolditalic]
\setmathfont{texgyrepagella-math.otf}
\setsansfont{texgyreheros}[
  Extension      = .otf,
  UprightFont    = *-regular,
  BoldFont       = *-bold,
  ItalicFont     = *-italic,
  BoldItalicFont = *-bolditalic,
  Scale          = MatchLowercase]
\setmonofont{FiraMono-Regular.otf}[
  BoldFont       = FiraMono-Bold.otf,
  Scale          = MatchLowercase]

% ---------- page geometry --------------------------------------------
% "8-page-sheet" (8 \textit{na yok}) standard textbook size: 19 x 26.5 cm
\usepackage[
  paperwidth  = 190mm,
  paperheight = 265mm,
  inner       = 26mm,
  outer       = 20mm,
  top         = 24mm,
  bottom      = 26mm,
  headheight  = 22pt,
  headsep     = 10mm,
  footskip    = 12mm]{geometry}

\usepackage{microtype}
\usepackage{setspace}
\setstretch{1.08}
\usepackage[skip=0.55em plus 0.1em]{parskip}
\setlength{\parindent}{0pt}
\usepackage{etoolbox}

% ---------- maths, tables, graphics ----------------------------------
\usepackage{amsmath,mathtools}
\usepackage{booktabs,array,multirow,tabularx,makecell}
\usepackage{graphicx}
\usepackage[export]{adjustbox}
\usepackage{xcolor}
\usepackage{tikz}
\usetikzlibrary{arrows.meta,positioning,calc,shapes.geometric,fit,automata}
\usepackage{pgfplots}
\pgfplotsset{compat=1.18}
\usepackage[siunitx]{circuitikz}
\usepackage{siunitx}
\sisetup{group-separator={,}}
\graphicspath{{figures/}{images/}{../public/}}

% ---------- colours ---------------------------------------------------
\definecolor{kmutnborange}{HTML}{E4610F}
\definecolor{kmutnbgreen}{HTML}{0B8A3E}
\definecolor{ink}{HTML}{1F2937}
\definecolor{accent}{HTML}{9A3F0C}
\definecolor{accentlight}{HTML}{FBEBDD}
\definecolor{bluelight}{HTML}{E8EFF8}
\definecolor{blueline}{HTML}{2B5C9E}
\definecolor{greenlight}{HTML}{E6F4EA}
\definecolor{greyline}{HTML}{8A8F98}
\definecolor{codebg}{HTML}{F6F7F9}
\definecolor{codekw}{HTML}{1D4E9E}
\definecolor{codecm}{HTML}{4C7A34}
\definecolor{codestr}{HTML}{A03A1B}
\definecolor{linkblue}{HTML}{1B4F9C}

% ---------- lists -----------------------------------------------------
\usepackage[inline]{enumitem}
\setlist{nosep, topsep=0.3em, itemsep=0.15em, leftmargin=1.5em}
\newlist{exlist}{enumerate}{2}
\setlist[exlist,1]{label=\textbf{\thechapter.\arabic*}, leftmargin=2.6em, itemsep=0.5em, topsep=0.4em, nosep=false}
\setlist[exlist,2]{label=(\alph*), leftmargin=1.8em, itemsep=0.15em}

% ---------- boxes -----------------------------------------------------
\usepackage[most]{tcolorbox}
\tcbset{
  enhanced, breakable, boxrule=0.6pt, arc=2pt, left=8pt, right=8pt, top=6pt, bottom=6pt,
  fonttitle=\sffamily\bfseries\small, before skip=0.8em, after skip=0.8em}

% Learning objectives (opening of every chapter)
\newtcolorbox{objectives}{
  colback=accentlight, colframe=kmutnborange, boxrule=0pt, leftrule=3pt, arc=0pt,
  title={Learning Objectives}, coltitle=accent, colbacktitle=accentlight,
  attach boxed title to top left={yshift=0pt}, boxed title style={boxrule=0pt, arc=0pt, left=0pt}}

% Worked examples (numbered within chapter)
\newtcbtheorem[number within=chapter]{example}{Example}{
  colback=white, colframe=blueline, colbacktitle=bluelight, coltitle=blueline,
  fonttitle=\sffamily\bfseries\small, attach boxed title to top left={yshift=-\tcboxedtitleheight/2},
  boxed title style={boxrule=0.6pt, colframe=blueline, arc=2pt}, top=10pt}{ex}

% Definitions / key ideas
\newtcbtheorem[number within=chapter]{definition}{Definition}{
  colback=greenlight, colframe=kmutnbgreen, colbacktitle=kmutnbgreen, coltitle=white,
  fonttitle=\sffamily\bfseries\small, boxrule=0.6pt, top=10pt,
  attach boxed title to top left={yshift=-\tcboxedtitleheight/2},
  boxed title style={arc=2pt}}{def}

% Key points / design rules
\newtcolorbox{keypoint}[1][Key Point]{
  colback=bluelight, colframe=blueline, boxrule=0pt, leftrule=3pt, arc=0pt,
  title={#1}, coltitle=blueline, colbacktitle=bluelight,
  attach boxed title to top left={yshift=0pt}, boxed title style={boxrule=0pt, arc=0pt, left=0pt}}

% Chapter summary
\newtcolorbox{chaptersummary}{
  colback=codebg, colframe=greyline, title={Chapter Summary}, coltitle=ink, colbacktitle=codebg,
  boxrule=0.6pt, attach boxed title to top left={yshift=-\tcboxedtitleheight/2},
  boxed title style={boxrule=0.6pt, colframe=greyline, arc=2pt}, top=10pt}

% Caution / common mistakes
\newtcolorbox{caution}[1][Common Mistake]{
  colback=white, colframe=accent, boxrule=0pt, leftrule=3pt, arc=0pt,
  title={#1}, coltitle=accent, colbacktitle=white,
  attach boxed title to top left={yshift=0pt}, boxed title style={boxrule=0pt, arc=0pt, left=0pt}}

% ---------- captions, floats ------------------------------------------
\usepackage{float}
\usepackage[font={small}, labelfont={sf,bf,color=accent}, labelsep=period, format=hang,
            justification=raggedright, singlelinecheck=true, skip=6pt]{caption}
\usepackage{subcaption}
\captionsetup[table]{position=top}

\usepackage{rotating}
% \figw{<width>}{<file>}{<caption>}{<label>} : figure with an explicit width (set by the layout script)
\newcommand{\figw}[4]{\begin{figure}[htbp]\centering\includegraphics[width=#1]{#2}\caption{#3}\label{#4}\end{figure}}
\newcommand{\figside}[4]{\begin{sidewaysfigure}[p]\centering\includegraphics[width=#1]{#2}\caption{#3}\label{#4}\end{sidewaysfigure}}
% \wavesplit{<name>}{<n>} : a wide simulation screenshot cut into n strips (each keeps the signal-name column) so that the text stays readable
\newcommand{\wavesplit}[2]{\foreach \wsi in {1,...,#2}{\includegraphics[width=\linewidth]{#1_p\wsi}\par\vspace{1.5pt}}}
% \fig[<graphicx keys>]{<file>}{<caption>}{<label>}   -- single figure
\newcommand{\fig}[4][]{%
  \begin{figure}[htbp]
    \centering
    \includegraphics[max width=\linewidth, max height=0.42\textheight, #1]{#2}
    \caption{#3}\label{#4}
  \end{figure}}

% ---------- listings (VHDL) -------------------------------------------
\usepackage{listings}
\lstdefinelanguage{VHDLx}{
  language=VHDL, morekeywords={%
    library,use,entity,architecture,is,of,begin,end,port,generic,map,in,out,inout,buffer,
    signal,variable,constant,process,if,then,elsif,else,case,when,others,for,loop,while,
    component,package,body,type,subtype,array,downto,to,wait,until,after,report,assert,
    rising_edge,falling_edge,with,select,generate,and,or,not,nand,nor,xor,xnor,mod,rem,abs},
  sensitive=false, morecomment=[l]{--}}
\lstdefinestyle{vhdl}{
  language=VHDLx, basicstyle=\ttfamily\footnotesize, keywordstyle=\color{codekw}\bfseries,
  commentstyle=\color{codecm}, stringstyle=\color{codestr},
  backgroundcolor=\color{codebg}, frame=single, rulecolor=\color{greyline!50},
  numbers=left, numberstyle=\tiny\color{greyline}, numbersep=6pt, xleftmargin=1.6em,
  framexleftmargin=1.6em, breaklines=true, breakatwhitespace=false, tabsize=2,
  showstringspaces=false, columns=fullflexible, keepspaces=true, aboveskip=0.6em, belowskip=0.6em,
  captionpos=b, upquote=true, literate={↔}{$\leftrightarrow$}1}
\lstset{style=vhdl}
\renewcommand{\lstlistingname}{Listing}
\lstnewenvironment{vhdl}[1][]{\lstset{style=vhdl,#1}}{}
% \vhdlfile[<options>]{<path>}{<caption>}{<label>}
\newcommand{\vhdlfile}[4][]{%
  \lstinputlisting[style=vhdl, caption={#3}, label={#4}, #1]{#2}}

\newcommand{\vhdlfilesmall}[4][]{%
  \lstinputlisting[style=vhdl, basicstyle=\ttfamily\scriptsize, caption={#3}, label={#4}, #1]{#2}}

% inline code
\DeclareRobustCommand{\code}[1]{\texttt{\detokenize{#1}}}

% ---------- index, glossary -------------------------------------------
\usepackage[noautomatic]{imakeidx}
\makeindex[intoc, title={Index}, columns=2, options={-s indexstyle.ist}]
% \term{<phrase>}  : emphasise first use AND add to the index
\newcommand{\term}[1]{\textit{#1}\index{#1}}
% \termx{<shown text>}{<index entry>} : same, with a different index entry
\newcommand{\termx}[2]{\textit{#1}\index{#2}}

\newenvironment{glossarylist}{%
  \begin{description}[font=\sffamily\bfseries, leftmargin=0pt, labelindent=0pt,
                      itemsep=0.5em, parsep=0.2em, style=unboxed]}{\end{description}}
\newcommand{\glossaryentry}[2]{\item[#1] #2}

% ---------- headers, footers, headings --------------------------------
\usepackage{fancyhdr}
\usepackage{titlesec}

\pagestyle{fancy}
\fancyhf{}
\renewcommand{\headrulewidth}{0.4pt}
\renewcommand{\headrule}{\hbox to\headwidth{\color{greyline}\leaders\hrule height \headrulewidth\hfill}}
\fancyhead[LE]{\sffamily\small\thepage\quad\textcolor{greyline}{|}\quad\nouppercase{\leftmark}}
\fancyhead[RO]{\sffamily\small\nouppercase{\rightmark}\quad\textcolor{greyline}{|}\quad\thepage}
\renewcommand{\chaptermark}[1]{\markboth{\ifnum\value{chapter}>0 \thechapter\ \fi#1}{}}
\renewcommand{\sectionmark}[1]{\markright{\ifnum\value{chapter}>0 \thesection\ \fi#1}}
\fancypagestyle{plain}{\fancyhf{}\renewcommand{\headrulewidth}{0pt}\fancyfoot[C]{\sffamily\small\thepage}}

\titleformat{\chapter}[display]
  {\sffamily\bfseries}
  {\color{kmutnborange}\fontsize{54}{54}\selectfont\thechapter\vspace{4pt}\hrule height 1.6pt\vspace{10pt}}
  {0pt}
  {\color{ink}\fontsize{26}{30}\selectfont}
\titlespacing*{\chapter}{0pt}{-6pt}{22pt}
\titleformat{name=\chapter,numberless}[display]
  {\sffamily\bfseries}
  {}
  {0pt}
  {\color{ink}\fontsize{26}{30}\selectfont\vspace{-4pt}\hrule height 1.6pt\color{kmutnborange}\vspace{14pt}}
\titleformat{\section}{\sffamily\bfseries\Large\color{ink}}{\color{kmutnborange}\thesection}{0.7em}{}
\titleformat{\subsection}{\sffamily\bfseries\large\color{ink}}{\color{kmutnborange}\thesubsection}{0.7em}{}
\titleformat{\subsubsection}{\sffamily\bfseries\normalsize\color{ink}}{}{0pt}{}
\titlespacing*{\section}{0pt}{1.6em plus 0.3em minus 0.2em}{0.6em}
\titlespacing*{\subsection}{0pt}{1.2em plus 0.2em minus 0.2em}{0.4em}
\titlespacing*{\subsubsection}{0pt}{0.9em}{0.2em}

% Unnumbered "end-of-chapter" headings that still land in the TOC
\newcommand{\chapsec}[1]{\section*{#1}\addcontentsline{toc}{section}{#1}\markright{#1}}

% ---------- table of contents formatting ------------------------------

% ---------- cross references & hyperlinks -----------------------------
\usepackage[
  colorlinks=true, linkcolor=linkblue, citecolor=linkblue, urlcolor=linkblue,
  bookmarksnumbered=true, pdfborder={0 0 0},
  pdftitle={Digital Circuits and Microprocessor Fundamentals},
  pdfauthor={Ruslee Sutthaweekul},
  pdfsubject={Textbook for Courses 010113025 and 010153101, KMUTNB}]{hyperref}
\usepackage[capitalise,noabbrev]{cleveref}
\crefname{listing}{Listing}{Listings}
\crefname{tcb@cnt@example}{Example}{Examples}
\Crefname{tcb@cnt@example}{Example}{Examples}
\crefname{tcb@cnt@definition}{Definition}{Definitions}
\Crefname{tcb@cnt@definition}{Definition}{Definitions}
\Crefname{listing}{Listing}{Listings}

% ---------- bibliography (BibTeX + IEEEtran numeric style) -------------
\renewcommand{\bibname}{Bibliography}

% ---------- maths shorthands used throughout --------------------------
\newcommand{\NOT}[1]{\overline{#1}}
\newcommand{\bin}[1]{\ensuremath{\mathtt{#1}}}              % binary string, e.g. \bin{1001}
\newcommand{\minterm}[1]{\ensuremath{m_{#1}}}
\newcommand{\maxterm}[1]{\ensuremath{M_{#1}}}
\newcommand{\Sm}{\ensuremath{\sum m}}
\newcommand{\PM}{\ensuremath{\prod M}}
\newcommand{\xor}{\ensuremath{\oplus}}
\newcommand{\rtl}[2]{\ensuremath{#1 \leftarrow #2}}         % register transfer: R1 <- R2
\newcommand{\dc}{\ensuremath{\mathrm{d}}}                   % don't care in a K-map

% “Note to the reader” helper (used sparingly)
\newcommand{\hint}[1]{\par\smallskip\noindent\textsf{\textbf{\textcolor{accent}{Hint.}}}~#1}

% truly empty blank pages before openright chapters
\makeatletter
\def\cleardoublepage{\clearpage\if@twoside\ifodd\c@page\else\hbox{}\thispagestyle{empty}\newpage\fi\fi}
\makeatother
